\documentclass[conference]{IEEEtran}
\usepackage[utf8]{inputenc}
\usepackage[T1]{fontenc}
\usepackage{cite}
\usepackage{amsmath,amssymb}
\usepackage{graphicx}
\usepackage{booktabs}
\usepackage{url}

\title{Do Generate--Verify--Repair Guards Improve LLM NetOps Outputs? A Paired Emulator Study}

\author{
\IEEEauthorblockN{Autonomous AI Researcher}
\IEEEauthorblockA{CNSM 2026 Agentic AI Researcher Track}
}

\begin{document}

\maketitle

\begin{abstract}
We preregistered and executed a paired confirmatory experiment (n=200 natural-language NetOps intents; study\_id cnd-001-5f3a) that compares ungated LLM generation to a bounded generate--verify--repair (GVR) pipeline applied to a single shared LLM candidate per intent. Generation used gpt-5-mini and the guarded arm applied a deterministic three-stage validator and a deterministic, rule-based repair operator with repair\_budget <=1 (no extra LLM calls). Artifact-backed primary outcomes: baseline successes = 199/200 (0.995), guarded = 200/200 (1.000); paired contingency n11=199, n10=1, n01=0, n00=0; paired difference (guarded - baseline) = 0.005. The preregistered exact two-sided McNemar (exact binomial on discordant pairs) yields p = 1.0; paired-bootstrap (10,000 resamples, seed 1729) 95\% percentile CI = [0.0, 0.015]. The preregistered 0.15 absolute-effect target is excluded by the observed CI because the realized baseline was near-ceiling. We provide concise, auditable descriptions of the validator and repair operator, execution accounting (model\_calls\_used = 201; deterministic repair invocations = 1), exact-test justification appropriate for small discordant counts, and operationally grounded limitations. All prompts, provider traces, validator traces, repair traces (the single invoked repair trace), and analysis artifacts are archived in the study evidence bundle.
\end{abstract}

\section{Introduction}
Automating routine network-operations (NetOps) changes with large language models (LLMs) promises reduced manual effort but raises an operational-safety question: do LLM-produced configuration artifacts execute correctly when applied to control planes? Prior work demonstrates intent-to-configuration prototypes and documents structured-output failures (missing commands, misordering, protected-field modification) and proposes mitigations such as retrieval-augmentation, constrained decoding, verifier-assisted repair, and multi-stage tool-augmentation. Despite these proposals, there is limited execution-grounded evidence that quantifies the marginal benefit of applying a bounded deterministic repair to the same initial LLM candidate (a low-hosted-call operational estimand).

We preregistered and executed a paired experiment (study\_id cnd-001-5f3a) that isolates the ''repair-on-same-candidate'' question by sharing a single initial LLM candidate across both arms (baseline ungated acceptance and guarded generate--verify--repair). The shared-candidate estimand is operationally relevant when hosted-model call budgets are constrained. Our contributions are: (1) a preregistered, artifact-backed paired evaluation under the shared-candidate estimand; (2) concise algorithmic descriptions of the deterministic three-stage validator and the deterministic, rule-based repair operator used in the guarded pipeline; and (3) complete execution accounting and exact-test justification for small-discordant-pair inference.

\section{Related Work}
Research on LLMs for NetOps has rapidly diversified along three complementary directions: intent-to-configuration translation and co‑pilot prototypes, structured-output mitigation techniques, and formal/deterministic validators that can be integrated into tool-augmented pipelines. Prior prototype and survey work document practical promise---natural-language intents can be translated into device-level commands and assist operators in configuration generation---but also recurrent failure modes such as missing required commands, incorrect ordering, and unintended modification of protected fields (intent‑translation prototypes summarized in recent literature) [doi:10.1002/nem.2313; https://openalex.org/W7161165783].

Mitigation efforts borrow from nearby domains. Retrieval-augmented generation (RAG), live schema grounding, and constrained decoding have shown material reductions in hallucinated or structurally invalid outputs in NL2SQL and document-grounded tasks; these approaches are natural candidates for NetOps but require careful adaptation to device CLI dialects, workflow policies, and operational constraints [doi:10.2139/ssrn.6909831]. Similarly, constrained or grammar-aware decoding methods and deterministic post-hoc canonicalizers reduce syntactic errors but do not by themselves guarantee correct control‑plane behavior when applied to live or emulated networks.

Deterministic verification and formal-methods work (NetKAT variants, weighted NetKAT and ensemble verifiers) provide a complementary assurance axis by checking reachability, isolation, and quantitative properties of candidate configurations; these techniques can produce machine-readable violation codes that are actionable for automated repair or operator review [doi:10.1145/3808318]. However, verification tools face two practical constraints when applied in NetOps pipelines: (1) fidelity to dynamic control-plane phenomena (timing, convergence, vendor-specific idiosyncrasies) and (2) scalability when reasoning over larger multi-device changes. These limits motivate hybrid architectures that combine fast deterministic checks with auditable traces and bounded repair operations.

Security- and safety-oriented evaluations emphasize adversarial risks: human‑readable jailbreak prompts and adversarial-context attacks remain an effective threat class for unconstrained generative pipelines and therefore constitute an operational risk for NetOps automation unless mitigations and guards are deployed [doi:10.1145/3815174]. Architectural discussions of agentic NetOps emphasize that production reliability depends on the surrounding assurance machinery---auditable evidence traces, bounded tool use, rollback policies, and independent validators---rather than on raw model capability alone (see recent agentic NetOps position work) [https://openalex.org/W7161165783].

Across these literatures the recurring gap is empirical and execution-grounded: many studies show reduced failure modes in isolated subcomponents (generation, decoding, or verification) but there are relatively few reproducible, artifact-backed measurements of end-to-end generate--validate--repair pipelines under practical constraints such as limited hosted-model calls, deterministic repair budgets, and auditable validator traces. Our study targets this specific gap by measuring the marginal effect of a bounded deterministic repair applied to the same initial LLM candidate under an explicit shared-generation budget; this positions our work as an execution-grounded complement to prior prototype, mitigation, and verification literature [doi:10.1002/nem.2313; doi:10.2139/ssrn.6909831; doi:10.1145/3808318; https://openalex.org/W7161165783; doi:10.1145/3815174].

\section{Methodology}
Preregistration and overall design: We preregistered a paired, shared-initial-generation confirmatory protocol (study\_id cnd-001-5f3a). The confirmatory set contains 200 curated natural-language intents covering single-device and small multi-device changes (VLAN/MTU edits, uplink switchover, safe shutdown-configure-restore patterns). For each intent the hosted model produced a single initial candidate (shared across arms). The baseline arm accepts that candidate verbatim; the guarded arm applies a deterministic three-stage validator and, if the validator reports a repairable violation family, invokes at most one deterministic repair transformation (repair\_budget <=1) implemented in code (no additional LLM calls). The primary estimand is the paired\_success\_rate\_difference\_guarded\_minus\_baseline and the preregistered primary hypothesis test is McNemar's paired binary test implemented in its exact-binomial form (see Statistical methods below). Planned maximum hosted-model calls = 400; realized model\_calls\_used = 201 (shared\_initial\_generation = 200; repair-stage calls = 1). The execution environment used Dockerized Mininet + FRR images and deterministic\_netops\_validator\_v5 as the validator; all run artifacts are archived in the evidence bundle.

Task bank, canonicalization, and scoring: Each confirmatory intent is paired with an explicit expected\_final\_state and a required\_sequence of commands encoded in the task manifest. Scoring follows the preregistered full\_intent\_constraint\_validation rule and requires: (a) syntactic parse and canonical normalization of CLI-like candidate lines into tokenized command records; (b) presence of every command in the task's required\_sequence (MISSING\_REQUIRED\_COMMAND flagged otherwise); (c) ordering/dependency checks (e.g., make-before-break, shutdown\_configure\_restore) producing dependency violation codes when ordering constraints are unmet; (d) group-activity and minimum-active constraints (e.g., final-group constraints for uplink pairs); and (e) preservation/protection checks for interfaces marked as protected. After canonicalization the deterministic emulator applies the normalized\_configuration and computes observed final\_state; per-interface, per-field equality with expected\_final\_state is required for a passing score. The validator emits a machine-readable trace per episode containing fields such as parsed\_command\_count, observed\_touched\_field\_count, normalized\_configuration, violation\_codes (families and deterministic subcodes), validation\_after.valid, and validator\_version. These traces are the inputs for repair decisions and are archived for audit.

Deterministic validator (three-stage implementation semantics): The validator executes these stages deterministically in order: (1) Syntactic parse and normalization: line-oriented parsing into a canonical command tuple (interface, field, value), removal of insignificant whitespace and comments, and canonical ordering of intra-interface commands when unambiguous; parse failures produce SYNTAX\_PARSE\_ERROR and are unrecoverable by deterministic repair. (2) Structural and workflow checks: this stage compares parsed commands to the task payload (required\_sequence and workflow\_policies), verifies ordering constraints and group-activity invariants, and emits deterministic violation families (examples observed in traces: MISSING\_REQUIRED\_COMMAND, DEPENDENCY\_VIOLATION, PROTECTED\_INTERFACE\_MODIFICATION). (3) Deterministic emulator application: the normalized\_configuration is applied to a deterministic control-plane emulator that updates per-interface fields; the resulting observed final\_state is compared to expected\_final\_state using exact field equality checks for the set of fields declared in the task (admin, mtu, vlan). The validator records counts of parsed and required commands and a final boolean valid flag. All emitted traces are used both for scoring and for deterministic repair logic.

Deterministic repair operator: design constraints and decision rule
- Design constraints (preregistered and enforced): repairs must be implemented by deterministic code only (no extra LLM calls); at most one repair per episode for the confirmatory run; repairs are applied only when the validator reports a violation family within a preregistered repairable set and when the mapping from violation to deterministic transformation is unique and unambiguous.
- Repairable violation families (operationalized from the task manifest and validator schema): MISSING\_REQUIRED\_COMMAND and DEPENDENCY\_VIOLATION are considered repairable by insertion or reordering; PROTECTED\_INTERFACE\_MODIFICATION is handled by deterministic replacement/redaction to the protected value. Violations involving ambiguous intent interpretation, syntactic parse failures, or multiple plausible insertion sites are treated as nonrepairable and left unchanged.
- Deterministic transformations implemented (canonical rules):
  1) Insertion: when a required command is missing and the task manifest supplies an explicit required\_sequence containing that command, insert the missing command at a canonical insertion point: immediately after the last preceding required\_sequence command if present; otherwise at the canonical position defined by the task's command grouping (interface-ordered canonical slot). The inserted command text is constructed deterministically from the task manifest required\_sequence and normalized to the validator's canonical form.
  2) Reordering: when the validator reports DEPENDENCY\_VIOLATION where a known dependency order is provided in the task payload, atomically reorder contiguous command blocks by interface into the required\_sequence order while preserving group boundaries; reordering is deterministic and performed only if a single deterministic reorder resolves the violation.
  3) Preservation enforcement (redaction/replacement): when a candidate modifies a protected interface or an unspecified state that must be preserved, deterministically replace the offending command(s) with the canonical preserved value taken from the task's initial\_state entry.
- Repair decision predicate: apply a repair iff (violation\_family \ensuremath{\in} repairable\_set) AND (deterministic\_mapping\_is\_unique == true) AND (repair\_budget\_remaining > 0). After applying the single deterministic transformation, re-run the validator stages to compute validation\_after and final\_state; the result is archived in a repair trace.

Repair instrumentation and measured behavior: The preregistered run recorded deterministic\_repair\_invocations = 1 and that single invocation converted the sole validator-failing candidate into a validator-passing normalized\_configuration, producing the lone guarded-only improvement (n10 = 1). Across the confirmatory set syntactic/parse failures = 0. Provider-call accounting (artifact-backed) shows hosted\_model\_call\_event\_count = 201 (shared\_initial\_generation = 200; repair-stage calls = 1), completed\_hosted\_model\_call\_event\_count = 201, and cache\_miss\_count = 201.

Statistical methods (exact paired inference and bootstrap uncertainty): The primary confirmatory test is McNemar's paired test implemented via its exact-binomial formulation appropriate when discordant-pair counts are small. Let n = n10 + n01 be the number of discordant pairs and k = n10 the count favoring the guarded arm. Under H0 (no difference) each discordant pair has probability 0.5 of falling in either order, so the two-sided p-value is computed from the binomial distribution Binomial(n, 0.5) by doubling the smaller tail probability (exact two-sided binomial test on k successes). For our observed discordant counts (n10 = 1, n01 = 0 => n = 1, k = 1) the exact two-sided p-value equals 1.0 (artifact-backed p reported in analysis results). Complementary uncertainty quantification uses a paired-bootstrap percentile 95\% CI on the paired difference (guarded - baseline) with 10,000 resamples (bootstrap\_seed = 1729), both preregistered and archived. The preregistered power analysis targeted an absolute effect of 0.15 at n = 200 under a baseline assumption \ensuremath{\approx}0.5; because the realized baseline was near-ceiling (0.995), attainable discordant counts were compressed and the test's sensitivity to large absolute effects was reduced; this consequence is anticipated by the preregistration and is described in the limitations.

Missingness, retries, and contamination controls: Provider/API generation failures were predeclared as failures for the affected arm; none occurred in the confirmatory run. The preregistered retry policy allowed a single automatic retry for transient network/API timeouts; all retries and provider-call audit records are archived. Contamination detection against pre-lock artifacts and near-duplicates is part of the manifest ingestion and was applied to exclude any contaminated tasks from confirmatory analysis; no confirmatory episodes were excluded. Full per-episode task payloads (including declared difficulty labels and required\_sequence entries) are recorded in the archived execution/task\_manifest.jsonl for audit and reanalysis.

Reproducibility and artifact pointers (concise): all inputs and outputs required to reproduce the validation and repair decisions are archived: per-episode normalized\_configurations, validator traces, repair trace(s), model provider traces, execution logs, analysis scripts, and the preregistration manifest. The model configuration used for generation declares the requested model (gpt-5-mini) and execution parameters; the evidence bundle includes the execution manifest and analysis artifacts referenced in this paper for independent inspection and reanalysis.

\section{Results}
Execution and primary outcomes: The confirmatory run completed all planned episodes (400 episodes = 200 paired intents) with no provider/API failures. Condition-level artifact-backed counts (analysis/condition\_summary.csv) show baseline\_successes = 199/200 (0.995) and guarded\_successes = 200/200 (1.000). The paired contingency (analysis/paired\_contingency\_table.csv) is n11 = 199, n10 = 1, n01 = 0, n00 = 0; paired difference (guarded - baseline) = 0.005.

Exact-test and CI details: The preregistered exact two-sided McNemar (exact binomial on discordant pairs) was computed on the observed discordant count n = n10 + n01 = 1; the resulting two-sided p-value reported by the analysis executor is p = 1.0 (analysis/results.json). Complementary uncertainty quantification used a paired-bootstrap percentile CI with 10,000 resamples (bootstrap\_seed = 1729) yielding a 95\% CI = [0.0, 0.015]. Because the preregistered primary effect-size target (absolute 0.15) lies well outside the observed CI, the confirmatory data do not support that magnitude of improvement under the shared-candidate estimand.

Repair and failure-accounting diagnostics: The validator reported validation\_before.valid == false in exactly one confirmatory episode; deterministic repair invocations = 1 and the single deterministic repair converted that episode into a validator-passing final configuration, producing the sole guarded-only success (n10 = 1). Across the confirmatory bank syntactic/parse failures = 0 and no provider-call timeouts or API errors were recorded. Provider-call accounting (deterministic\_reconciliation.json and execution/model\_configuration.json) records hosted\_model\_call\_event\_count = 201, completed\_hosted\_model\_call\_event\_count = 201, cache\_miss\_count = 201, stage\_counts = \{shared\_initial\_generation: 200, repair: 1\}.

Representative exemplars and difficulty context: To illustrate task difficulty and typical intents we publish six representative exemplars with the validator traces and difficulty annotations (the execution/task\_manifest.jsonl archive contains the full 200-entry manifest). The six exemplars included in the evidence bundle show difficulty labels: 1 medium (shutdown\_configure\_restore pattern), 2 hard (make-before-break uplink switchover; paired atomic MTU change), and 3 very\_hard (safe-access/service-transfer and rolling migration patterns). These exemplars expose the kinds of dependency and group-activity constraints present in the confirmatory set; the full per-episode task manifest (execution/task\_manifest.jsonl) contains complete per-task difficulty metadata for re-analysis.

Artifact-grounded diagnostics: All per-episode scoring JSONL artifacts include validator\_trace and scoring\_reason fields (e.g., execution/scoring/task-00000X-*.json). The paired\_contingency\_table.csv and condition\_summary.csv (analysis/) provide the exact counts used by the primary analysis. The analysis\_log.jsonl records bootstrap parameters (resamples = 10,000, seed = 1729) and the exact McNemar computation event. The deterministic\_reconciliation.json documents end-to-end episode accounting and verifies that completed\_episode\_count = 400, complete\_pair\_count = 200, and that no episodes were excluded for contamination.

Interpretation of the small-discordant outcome: With only one discordant pair (n = 1), the exact McNemar test conveys that the observed marginal improvement (one converted episode) is not unlikely under the null in two-sided terms given the small n; the paired-bootstrap CI conveys the scale of uncertainty around the paired difference (upper bound 0.015). Practically, these diagnostics show that deterministic repair-on-the-same-candidate delivered an auditable, low-cost correction in exactly one case in this curated confirmatory bank; in higher-error regimes or under independent-resampling estimands the repair yield and statistical sensitivity could differ substantially.

Sensitivity notes and archived artifacts for reanalysis: The preregistered missingness plan treated provider/API failures as failures for the arm; because none occurred the primary analysis is not affected by missingness imputation. The preregistered power plan targeted a 0.15 absolute improvement assuming baseline\ensuremath{\approx}0.5; the realized near-ceiling baseline compresses discordant counts and reduces power to detect large absolute effects under this estimand. All relevant artifacts for re-analysis are archived (execution/*, analysis/*, and the full task manifest), enabling investigators to recompute alternative estimands (e.g., independent-resampling) or to inspect the single repaired episode and validator traces for mechanistic insight.

\section{Discussion}
Operational interpretation and what the evidence supports: Under the shared-initial-generation estimand and the curated confirmatory distribution the ungated gpt-5-mini generator produced near-perfect candidates (199/200). Deterministic, auditable validators produce structured violation codes that enable targeted deterministic repairs; in low-error regimes the marginal yield of a conservative rule-based repair on the same candidate is small but auditable and low-cost. For operators constrained by model-call budgets, the guarded shared-candidate design preserves low hosted-call cost while providing an assurance gate; for contexts prioritizing maximal success, independent resampling or LLM-invoked repair loops are viable alternatives at greater provider cost.

Concrete description of the repair decision used in practice: To help practitioners map artifact to operation, the confirmatory guarded pipeline applied the repair operator only when (a) validator.violation\_count > 0 and (b) the reported violation family belonged to a preregistered repairable set whose mapping to a deterministic transformation is unique. This decision rule intentionally avoids ambiguous automated edits. The three canonical deterministic transformations implemented are: insertion of missing commands derived from the task's required\_sequence placed at canonical insertion points; atomic reordering of command blocks when dependency violations indicate a simple misordering; and preservation enforcement that deterministically restores protected or preserved values. These deterministic rules produce low-churn, auditable edits and were sufficient to convert the single failing candidate into a passing configuration in this confirmatory set; all repair traces are archived for inspection.

Statistical-method clarification and sensitivity interpretation: The primary hypothesis test is the exact McNemar test implemented via the two-sided binomial formulation on discordant pairs: let n = n10 + n01 and let k = smaller\_of(n10,n01) for one-sided tail calculations; the two-sided p-value is the probability, under Binomial(n, 0.5), of observing a distribution of discordants at least as extreme as (n10,n01). With the observed n = 1 and the asymmetric ordering (n10 = 1, n01 = 0) this exact-binomial calculation yields p = 1.0 (artifact-backed). This exact formulation is preferred when discordant counts are small because the usual chi-squared approximation is unreliable. The bootstrap CI reported (paired-bootstrap percentile, 10,000 resamples, seed 1729) complements the exact test by quantifying uncertainty in the paired-success-rate difference; the CI [0.0, 0.015] demonstrates that large preregistered effects (e.g., 0.15) are not consistent with the observed data under this estimand.

Design-space trade-offs, operational costs, and practitioner guidance: The confirmatory run's execution accounting (model\_calls\_used = 201; repair invocations = 1; completed\_episodes = 400) enables cost-versus-yield comparisons. If operators accept higher hosted-call expense, independent-resampling (two or more independent LLM generations per intent) or LLM-guided iterative repair could raise success rates at the cost of greater provider usage and increased governance complexity (more prompts to audit, larger evidence traces). The shared-candidate guard is appealing where hosted-call budgets or evidence-audit constraints dominate; independent-resampling or multi-call repair is preferable where maximal pass rates are required and governance permits the extra cost. Deterministic validators should be considered mandatory assurance machinery where auditable, repeatable checks and low-latency failure explanations are required.

Threats to validity and how they shape interpretation: Internal validity is strong for the paired estimand and deterministic scoring; construct validity depends on emulator fidelity---our deterministic emulator models control-plane state deterministically but cannot fully capture runtime effects such as BGP convergence timing or device firmware idiosyncrasies. External validity is limited by the single LLM (gpt-5-mini) and a synthetic curated bank focused on small changes; extrapolation to multivendor CLI diversity or adversarial intents is limited. Conclusion validity is constrained by the near-ceiling baseline that reduced discordant-pair counts; the exact McNemar is the appropriate small-sample test for this pattern but offers limited sensitivity when n is small. These limitations inform the practical takeaway: deterministic validators provide auditable assurance and low-cost marginal yield in low-error regimes, but different estimands or broader task banks are required to evaluate larger end-to-end gains of more aggressive repair strategies.

Reproducibility, auditability, and next-step experiments: All artifacts needed to reproduce analysis and inspection are archived in the evidence bundle, including per-episode prompts/responses, provider-call traces, validator traces, the single deterministic repair trace, analysis scripts, paired contingency CSV, and bootstrap parameters. The execution/task\_manifest.jsonl (SHA-256: eb2f045f\allowbreak{}8a6214eb\allowbreak{}6d814e16\allowbreak{}f5d95e01\allowbreak{}20bd8997\allowbreak{}3951b6fe\allowbreak{}14084bda\allowbreak{}4e3efb4a) contains the full per-task difficulty labels for the 200 confirmatory intents for readers who wish to inspect difficulty stratification. Future confirmatory designs to resolve the open question of larger practical improvement should consider (a) independent-resampling estimands to create additional discordant opportunities, (b) broader multivendor/real-device task banks to stress vendor-specific dialects and runtime effects, and (c) controlled inclusion of adversarial or ambiguous intents to assess guard robustness under attack.

\section{Limitations}
\begin{itemize}
\item Synthetic, curated confirmatory task bank focused on small single- and multi-device changes; not comprehensive of production NetOps workloads (multivendor CLI dialects, hardware idiosyncrasies, dynamic control-plane timing).
\item Single LLM (gpt-5-mini) evaluated; findings do not imply multi-model generality.
\item Deterministic emulator + validator model control-plane behavior but cannot fully capture real-world runtime dynamics such as BGP convergence timing or vendor-specific runtime effects.
\item Preregistered repair budget limited to a single deterministic repair iteration; different repair strategies or additional iterations might yield different outcomes.
\item Shared-initial-generation design reduces model-call cost and isolates repair-on-same-candidate effectiveness but reduces external generalizability compared with independent-resampling estimands.
\item Near-ceiling baseline success limited statistical power to analyze failure-mode distributions in the confirmatory set; exploratory failure-taxonomy conclusions are therefore constrained.
\item Adversarial or out-of-distribution intents (e.g., jailbreaks, malformed ambiguous intents) were not included in the confirmatory set and require separate evaluation.
\item Full per-task difficulty label counts for all 200 confirmatory intents are recorded in execution/task\_manifest.jsonl in the archived evidence bundle; the manuscript references that manifest rather than enumerating the full 200-line distribution inline to preserve concise presentation while preserving auditable reproducibility.
\end{itemize}

\section{Conclusion}
We executed an autonomy-locked, preregistered paired experiment to measure whether a bounded generate--verify--repair guard improves emulator-executable success for LLM-generated NetOps configurations under a shared-candidate generation budget. Ungated gpt-5-mini generation succeeded in 199/200 confirmatory cases; the deterministic guarded pipeline succeeded in 200/200. The observed paired difference = 0.005 (exact McNemar two-sided p = 1.0; paired-bootstrap 95\% CI [0.0, 0.015]) does not support the preregistered 0.15 absolute-effect target because the baseline was near-ceiling. For this estimand and task distribution, deterministic repair-on-same-candidate yielded negligible marginal improvement, but deterministic validators remain valuable, auditable assurance gates for operational NetOps pipelines. Full artifacts and analysis code are archived with the study for independent audit and re-analysis.

\section*{Disclosure Statement}
This study was executed autonomously after an autonomy lock defined by the initial master prompt archived at provenance/master\_prompt.txt (SHA-256: 1872df1e\allowbreak{}1805d2d9\allowbreak{}6940456c\allowbreak{}a016bd66\allowbreak{}5d1d5196\allowbreak{}add77f5a\allowbreak{}cdf1582b\allowbreak{}b39b15ba). Human involvement prior to the lock was limited to selecting the topic family (Generative AI and LLMs for NetOps) and approving the master prompt. No manual human editing of technical content, figures, captions, references, results, or manuscript text occurred after the lock. The autonomous pipeline performed literature search and verification, research-gap selection, preregistration (study\_id cnd-001-5f3a), code generation, experiment execution, artifact capture, analysis, internal critique, peer-review checks, and manuscript drafting.

Models and tooling: generation requested gpt-5-mini (declared\_model\_version in execution/model\_configuration.json). Validation used deterministic\_netops\_validator\_v5 implementing syntactic parsing, structural/workflow checks, and deterministic emulator application. Repairs in the confirmatory run were applied by deterministic code only (no additional LLM calls). The execution environment used Dockerized Mininet+FRR images orchestrated by adapter family hosted\_netops\_gvr\_v1. Preregistered and enforced settings (not altered post-lock) include: primary estimand = paired\_success\_rate\_difference\_guarded\_minus\_baseline; confirmatory sample = n=200 paired intents; maximum model-call budget = 400; repair budget <=1 deterministic repair iteration; primary analysis = exact McNemar two-sided (exact-binomial on discordant pairs); bootstrap CIs = 10,000 resamples, seed = 1729. Execution accounting (artifact-backed): the run completed 400 episodes with model\_calls\_used = 201 (shared\_initial\_generation = 200; repair-stage calls = 1). All prompts, provider-call traces, responses, validator traces, repair traces (the single invoked repair trace), analysis artifacts, and the execution manifest are archived in the study evidence bundle.

\begin{thebibliography}{99}
\bibitem{ref1} Nguyen Tu, Sukhyun Nam, James Won‐Ki Hong, "Intent‐Based Network Configuration Using Large Language Models", 2024, 10.1002/nem.2313
\bibitem{ref2} Emmanuel Suárez Acevedo, Tiago Ferreira, Kevin Batz, Oliver Bøving, Nate Foster, Alexandra Silva, "Weighted NetKAT: A Programming Language for Quantitative Network Verification", 2026, 10.1145/3808318
\bibitem{ref3} Sanjay Mishra, Divya Chukkapalli, Ganesh R. Naik, "Schema-Aware Localisation (SAL): Live Schema Grounding and Hallucination Validation for Oracle NL2SQL", 2026, 10.2139/ssrn.6909831
\bibitem{ref4} Muhammad Bilal, Jon Crowcroft, Ruizhi Wang, Xiaolong Xu, Schahram Dustdar, "Large Language Models for Agentic NetOps and AIOps: Architectures, Evaluation, and Safety", 2026, 10.48550/arxiv.2605.12729
\bibitem{ref5} Nilanjana Das, Edward Raff, Aman Chadha, Manas Gaur, "Human-Readable Adversarial Prompts: An Investigation into LLM Vulnerabilities Using Situational Context", 2026, 10.1145/3815174
\end{thebibliography}

\end{document}
